\section*{Appendix: Worked Example --- 2017 Natural Catastrophe}

This appendix traces the exposure-adjustment formula through a single
historical event: the 2017 HIM (Harvey--Irma--Maria) hurricane season.
Two syndicates are chosen to illustrate how raw severity can differ
by a factor of seven while the standardised severity is identical.

\subsection*{Event and syndicates}

Event identifier: \texttt{2017\_natural\_cat} (30 syndicates observed).
Raw severity $S_{\text{raw},A}$ across the 30 syndicates ranges from
$-0.01\%$ to $+55.5\%$.

\medskip
\begin{tabular}{lrrl}
\toprule
 & Syndicate A (1414) & Syndicate B (4444) & \\
\midrule
Reserves $R_s$ (\pounds m) & 19.3 & 389.8 & \\
Property weight $w_{s,\text{Prop}}$ & 0.319 & 0.084 & \\
Reins-Property weight & 0.582 & 0.437 & \\
Professional Lines weight & 0.018 & 0.426 & \\
Other LoB weights & 0.081 & 0.053 & \\
\addlinespace
Property severity $s_{\text{Prop}}$ (uncapped) & 11.50 & 6.13 & \\
Property severity $s_{\text{Prop}}$ (capped at $\pm 5$) & 5.00 & 5.00 & \\
All other $s_\ell$ & 0.00 & 0.00 & \\
\addlinespace
\textbf{Raw severity} $S_{\text{raw},A}$ & \textbf{3.67\%} & \textbf{0.51\%} & ratio $= 7.2\times$ \\
\bottomrule
\end{tabular}

\medskip
Syndicate~1414 is a small, concentrated property/reinsurance-property
book.  The 2017 hurricanes produced an 11.5\% per-unit severity on its
Property exposure; because 32\% of the portfolio sits in Property,
this transmits as 3.67\% at the total-portfolio level.

Syndicate~4444 is 20$\times$ larger with only 8\% in Property.  Its
per-unit Property severity was 6.1\%, but because Property is a minor
slice of the portfolio the total hit was only 0.51\%.

\subsection*{Step 1: Mix standardisation}

Query portfolio (property-heavy):
$\mathbf{w}_q = (0.60,\; 0.20,\; 0.10,\; 0.10)$ across Property,
Casualty, Marine, Professional Lines.

\[
S_{\text{std}} = \mathbf{w}_q \cdot \mathbf{s}_{\text{lob}}
  = 0.60 \times 5.00 + 0.20 \times 0 + 0.10 \times 0 + 0.10 \times 0
  = 3.00\%
\]

Both syndicates converge to $S_{\text{std}} = 3.00\%$.  The 7$\times$
difference in raw severity was entirely attributable to LoB-mix
heterogeneity, not to a difference in the underlying event severity.

\subsection*{Step 2: Size adjustment}

For a \pounds 200m query portfolio with $\beta_{\text{weighted}} = -0.358$
and $R_{\text{ref}} = 500$:

\[
A = \left(\frac{200}{500}\right)^{-0.358} = 0.40^{-0.358} = 1.388
\]

\[
S_{\text{adj}} = S_{\text{std}} \times A = 3.00\% \times 1.388 = 4.16\%
\]

The size factor inflates the severity because the query portfolio
(\pounds 200m) is smaller than the reference (\pounds 500m): smaller
books experience proportionally larger reserve shocks for the same
event.

\subsection*{Summary}

\begin{tabular}{lrrr}
\toprule
 & Synd.\ 1414 & Synd.\ 4444 & Query portfolio \\
\midrule
Raw $S_{\text{raw},A}$ & 3.67\% & 0.51\% & --- \\
Mix-standardised $S_{\text{std}}$ & 3.00\% & 3.00\% & 3.00\% \\
Fully adjusted $S_{\text{adj}}$ & --- & --- & 4.16\% \\
\bottomrule
\end{tabular}

\medskip
The worked example demonstrates that the adjustment formula removes
idiosyncratic LoB-mix and size effects, producing a severity estimate
that answers the question: ``what would this event have meant for
\emph{my} portfolio?''
